% ============================================================
% APPENDIX A: Balance-of-Payments Decomposition and Central Bank
%             Solvency Ratio Formalization
% ============================================================

\section[Formalization of Solvency Ratio and Balance-of-Payments Accounting]{Formalization of the Central Bank Solvency Ratio and\\Balance-of-Payments Accounting}
\label{app:bop_levr}

\textbf{A.1} This appendix develops the formal balance-sheet accounting, the exact discrete law of motion, and the structural constraint-binding conditions governing the central bank \textbf{Solvency Ratio} ($\text{SolvR}_t$). Every accounting step is an exact discrete identity; no linear approximation or continuous-time distortion enters the derivation.

\subsection{The Solvency Ratio in Nominal Terms}
\label{app:solvr_nominal}

\textbf{A.2} All variables are denominated in current domestic currency (Chilean pesos/escudos, CLP) at time $t$, with the nominal exchange rate convention defined as $e_t$ (CLP per USD), such that a nominal depreciation corresponds to $e_t \uparrow$.

\textbf{A.3} The balance-sheet and external-account variables are defined as follows:
\begin{itemize}[leftmargin=2em]
	\item $M_t$: domestic monetary obligations (CLP), representing the liquid liabilities of the consolidated domestic financial system underwritten by the central bank.
	\item $IR_t$: stock of gross international reserves (USD), representing claims denominated in world money.
	\item $e_t$: nominal exchange rate (CLP per USD).
	\item $F_t \equiv e_t \cdot IR_t$: stock of world money held by the central bank, valued in domestic currency (CLP).
	\item $\text{BoP}_t \equiv \Delta F_t$: net balance-of-payments flow (CLP).
	\item $g_{M,t} \equiv \Delta M_t / M_t$: simple growth rate of domestic money obligations.
	\item $g_{F,t} \equiv \Delta F_t / F_t = \text{BoP}_t / (e_t \cdot IR_t)$: simple growth rate of the world-money asset base.
	\item $X_t$: export volume of primary and manufactured goods.
	\item $Q_t$: import volume of intermediate, consumption, and capital goods.
	\item $P_{Q,t}^{USD}$: import price index in foreign currency (USD).
	\item $\text{ToT}_t \equiv P_{X,t}^{USD} / P_{Q,t}^{USD}$: terms of trade.
	\item $D_t^{ext}$: stock of external debt (USD).
	\item $r_t$: average nominal interest rate on external debt.
	\item $FDI_t$: net foreign direct investment inflows (USD).
	\item $\Pi_t$: net foreign profit and dividend remittances (USD).
	\item $T_t$: net current external transfers (CLP).
\end{itemize}

Import volumes carry the symbol $Q_t$ throughout to prevent notation collisions with domestic money $M_t$. Discrete growth rates are evaluated on the contemporaneous base to maintain algebraic exactness.

\textbf{A.4} I define the central bank Solvency Ratio ($\text{SolvR}_t$) in nominal terms as:
\begin{equation}
	\label{app:eq:solvr_def}
	\text{SolvR}_t \;\equiv\; \frac{e_t \cdot IR_t}{M_t} \;=\; \frac{F_t}{M_t}
\end{equation}
The Solvency Ratio measures the stock of world-money assets held by the central bank per unit of unanchored domestic credit obligations underwritten by the monetary authority. Its algebraic inverse, $\text{LevR}_t \equiv 1/\text{SolvR}_t = M_t / (e_t \cdot IR_t)$, expresses the corresponding balance-sheet leverage ratio. I center the analytical framework on the Solvency Ratio because it directly formalizes central bank reserve backing and operationalizes the transposition of Minsky's financial postures to the peripheral central bank's own balance sheet \citep{Foley2003,ChamorroFutinico2025} (Section~\ref{sec:B5_foley_chamorro_minsky}).

\textbf{A.5} Taking continuous growth rates of Eq.~\eqref{app:eq:solvr_def} yields the continuous growth decomposition:
\begin{equation}
	\label{app:eq:solvr_growth}
	g_{\text{SolvR},t} \;=\; g_{F,t} \;-\; g_{M,t} \;=\; g_{e,t} \;+\; g_{IR,t} \;-\; g_{M,t}
\end{equation}
Equation~\eqref{app:eq:solvr_growth} isolates three distinct structural channels driving peripheral balance-sheet solvency:
\begin{enumerate}[leftmargin=2em]
	\item \textbf{External Balance-of-Payments Channel} ($g_{IR,t}$, expanding solvency): Export revenues and net foreign capital inflows accumulate foreign exchange reserves, bolstering central bank solvency. Current account deficits, foreign debt service, and transnational profit repatriation drain reserves, pulling $\text{SolvR}_t$ downward.
	\item \textbf{Valuation Channel} ($g_{e,t}$, expanding nominal solvency): Nominal currency depreciation ($e_t \uparrow$) increases the domestic-currency accounting valuation of foreign assets ($F_t = e_t \cdot IR_t$). While this mechanically cushions $\text{SolvR}_t$ in domestic currency units, it carries acute real macroeconomic costs by inflating imported intermediate input costs and triggering domestic wage-price defense spirals.
	\item \textbf{Domestic Dilution Channel} ($g_{M,t}$, compressing solvency): Broad money creation, fiscal monetization, and emergency Lender-of-Last-Resort liquidity expand domestic obligations from below, continuously diluting the world-money backing of the monetary base.
\end{enumerate}

\subsection{Current and Capital Account Decomposition}
\label{app:current_account}

\textbf{A.6} The evolution of the central bank's foreign-currency asset base is governed by the external balance of payments. In domestic currency terms, the current account ($CA_t$) and capital account ($KA_t$) are specified as:
\begin{equation}
	\label{app:eq:CA}
	CA_t \;=\; e_t \cdot P_{Q,t}^{USD}\!\left(\text{ToT}_t \cdot X_t - Q_t\right) \;-\; e_t \cdot r_t \cdot D_t^{ext} \;+\; T_t
\end{equation}
\begin{equation}
	\label{app:eq:KA}
	KA_t \;=\; e_t \cdot \Delta D_t^{ext} \;+\; e_t \cdot FDI_t \;-\; e_t \cdot \Pi_t
\end{equation}
The sum of the current and capital accounts determines the net balance-of-payments flow ($\text{BoP}_t$), which dictates the net change in world-money holdings:
\begin{equation}
	\label{app:eq:bop}
	\Delta F_t \;\equiv\; \Delta(e_t \cdot IR_t) \;=\; CA_t \;+\; KA_t \;\equiv\; \text{BoP}_t
\end{equation}

\subsection{Derivation of the Discrete Solvency Law of Motion}
\label{app:solvr_substitution}

\textbf{A.7} From Eq.~\eqref{app:eq:solvr_def}, the total stock of world-money assets held by the central bank decomposes into the product of the Solvency Ratio and domestic monetary obligations:
\begin{equation}
	\label{app:eq:solvr_sub}
	F_t \;=\; \text{SolvR}_t \cdot M_t
\end{equation}
Applying the exact discrete product rule $\Delta(A_t B_t) = B_t \Delta A_t + A_{t-1} \Delta B_t$ with $A_t = \text{SolvR}_t$ and $B_t = M_t$ yields:
\begin{equation}
	\label{app:eq:solvr_expand}
	\Delta F_t \;=\; M_t \cdot \Delta\text{SolvR}_t \;+\; \text{SolvR}_{t-1} \cdot \Delta M_t
\end{equation}
Dividing Eq.~\eqref{app:eq:solvr_expand} through by $M_t$ and substituting $\Delta F_t = \text{BoP}_t$ from Eq.~\eqref{app:eq:bop}:
\begin{equation}
	\frac{\text{BoP}_t}{M_t} \;=\; \Delta\text{SolvR}_t \;+\; \text{SolvR}_{t-1} \cdot \frac{\Delta M_t}{M_t}
\end{equation}
Rearranging terms isolates the exact discrete law of motion for the Solvency Ratio:
\begin{equation}
	\label{app:eq:solvr_lom}
	\boxed{\;\Delta\text{SolvR}_t \;=\; \frac{\text{BoP}_t}{M_t} \;-\; \text{SolvR}_{t-1} \cdot g_{M,t}\;}
\end{equation}
where $g_{M,t} \equiv \Delta M_t / M_t$ is the simple growth rate of domestic monetary liabilities.

\textbf{A.8} Equation~\eqref{app:eq:solvr_lom} delivers an intuitive and powerful economic decomposition:
\begin{itemize}[leftmargin=2em]
	\item The first term, $\text{BoP}_t / M_t$, represents the \textbf{external liquidity injection rate}---net foreign exchange inflows scaled directly by the existing domestic credit pyramid. Balance-of-payments surpluses expand the asset base and bolster solvency.
	\item The second term, $\text{SolvR}_{t-1} \cdot g_{M,t}$, represents the \textbf{domestic dilution drag}. Rapid domestic money growth ($g_{M,t} > 0$) dilutes the inherited backing ratio. The presence of the lagged state variable $\text{SolvR}_{t-1}$ is an exact algebraic consequence of the discrete product rule: domestic monetary expansion dilutes the preexisting stock of world-money cover.
\end{itemize}

\textbf{A.9} For completeness, the dual discrete law of motion for the leverage ratio ($\text{LevR}_t \equiv 1/\text{SolvR}_t$) follows symmetrically from $M_t = \text{LevR}_t \cdot F_t$:
\begin{equation}
	\label{app:eq:levr_lom}
	\Delta\text{LevR}_t \;=\; \frac{\Delta M_t}{e_t \cdot IR_t} \;-\; \text{LevR}_{t-1} \cdot g_{F,t}
\end{equation}
Both laws express the identical balance-sheet reality: in solvency form (Eq.~\ref{app:eq:solvr_lom}), foreign inflows drive recovery while domestic monetization acts as a dilution drag; in leverage form (Eq.~\ref{app:eq:levr_lom}), domestic money creation drives leverage expansion while reserve accumulation acts as an absorption multiplier.

\subsection{Analytical Convexity and Non-Linear Reserve Exhaustion}
\label{app:convexity}

\textbf{A.10} The structural vulnerability of peripheral central banks is amplified by the non-linear curvature of balance-sheet solvency. Taking the partial derivatives of $\text{SolvR}_t$ with respect to domestic obligations $M_t$ and foreign assets $F_t$:
\begin{equation}
	\label{app:eq:solvr_derivatives}
	\frac{\partial \text{SolvR}_t}{\partial M_t} \;=\; -\frac{F_t}{M_t^2} \;<\; 0, \qquad \frac{\partial^2 \text{SolvR}_t}{\partial M_t^2} \;=\; \frac{2 F_t}{M_t^3} \;>\; 0 \quad (\forall M_t > 0, F_t > 0)
\end{equation}
\begin{equation}
	\frac{\partial \text{SolvR}_t}{\partial F_t} \;=\; \frac{1}{M_t} \;>\; 0, \qquad \frac{\partial^2 \text{SolvR}_t}{\partial F_t^2} \;=\; 0
\end{equation}
As gross international reserves approach exhaustion ($IR_t \to 0 \implies F_t \to 0$), the absolute buffer $\text{SolvR}_t$ approaches zero. In this depleted state, any incremental withdrawal of foreign currency or external import shock triggers an acute percentage collapse in central bank credibility, permanently crippling the authority's capacity to anchor expectations.

\subsection{Constraint-Binding Conditions and Structural Minskian Regimes}
\label{app:constraint_cases}

\textbf{A.11} The central bank's solvency position deteriorates ($\Delta\text{SolvR}_t < 0$) whenever domestic monetary dilution outpaces net balance-of-payments inflows:
\begin{equation}
	\label{app:eq:binding}
	\Delta\text{SolvR}_t \;<\; 0 \;\iff\; \frac{\text{BoP}_t}{M_t} \;<\; \text{SolvR}_{t-1} \cdot g_{M,t} \;\iff\; g_{M,t} \;>\; \frac{\text{BoP}_t}{F_{t-1}}
\end{equation}
Even during an external commodity price boom ($\text{BoP}_t > 0$, such as the 1971--1972 copper price high), central bank solvency deteriorates if domestic credit expansion outpaces the growth rate of foreign reserves.

\textbf{A.12} In my structural Post-Keynesian framework, monetization is an endogenous institutional response to distributive conflict and foreign-exchange bottlenecks rather than an exogenous policy choice. The state-dependent parameter $\beta_{gM}(\text{SolvR}_{t-1})$ characterizes the regime-conditional absorption capacity of the peripheral monetary system. Following the 3-regime Threshold Vector Autoregression estimated in Section~\ref{sec:empirical_tournament}, the macro-system transitions across three discrete Minskian states:
\begin{enumerate}[leftmargin=2em]
	\item \textbf{Hedge Regime ($\text{SolvR}_{t-1} > \hat{\gamma}_H = 0.21$):} Reserve backing comfortably validates domestic liquid obligations. The central bank credibly maintains currency stability, and domestic monetary expansion is absorbed without price instability ($\beta_{\text{Hedge}} \approx 0.00\%$).
	\item \textbf{Speculative Regime ($0.14 < \text{SolvR}_{t-1} \le 0.21$):} Reserve buffers are partially depleted. Domestic credit growth begins to strain central bank credibility, and the pass-through of money to prices increases to $+0.54\%$.
	\item \textbf{Forced Ponzi Regime ($\text{SolvR}_{t-1} \le \hat{\gamma}_P = 0.14$):} Foreign exchange reserves are exhausted. The central bank is forced into emergency Lender-of-Last-Resort monetization to avert systemic enterprise insolvencies. Because world-money backing is gone, emergency liquidity cannot be absorbed by the balance sheet and passes directly into hyperinflation and currency flight ($\beta_{\text{Ponzi}} = +4.54\%$).
\end{enumerate}
This threshold formalization establishes that hyperinflation is a non-linear property of balance-sheet solvency exhaustion under external financial subordination, not a linear artifact of money creation.
